% ===================================================================
%  TABEL Nr. 3 — Lapsepõlv ja noorus.
%  Traduction 2026 d'apres texto/io, controlee sur texto/fr et
%  texto/fi. Consigne : voir texto/et/10-tabel-01.tex.
%
%  LE CALENDRIER DE L'ALINEA 1 DE LA SECONDE SCENE RENCONTRE LA
%  COLONNE ROMANCHE, ET LES DEUX LANGUES N'ONT AUCUN RAPPORT.
%
%  Le tableau 3 romanche avait releve que le romanche garde les mois
%  de Rome sauf deux, et que ces deux-la portent le nom du travail
%  qu'on y fait : « zercladur » le mois ou l'on bine, « fanadur » LE
%  MOIS OU L'ON FANE.
%
%  L'estonien a un second calendrier complet, encore vivant a cote
%  des noms latins, et il y nomme juillet EXACTEMENT DE MEME :
%
%    « heinakuu »   juillet, LE MOIS DU FOIN ;
%    « lõikuskuu »  aout, le mois de la moisson ;
%    « lehekuu »    mai, le mois des feuilles ;
%    « paastukuu »  mars, le mois du careme ;
%    « viinakuu »   octobre, le mois de l'eau-de-vie ;
%    « südakuu »    janvier, le mois du coeur de l'hiver ;
%  et six autres tires des saints du calendrier paysan :
%    « jürikuu » avril, « jaanikuu » juin, « mihklikuu » septembre,
%    « mardikuu » novembre, « jõulukuu » decembre, « küünlakuu »
%    fevrier.
%
%  DEUX LANGUES SANS CONTACT, UNE ROMANE ET UNE FINNO-OUGRIENNE,
%  NOMMENT JUILLET PAR LE FOIN. Ce n'est pas un emprunt, ce n'est pas
%  une parente : c'est que dans les deux pays juillet ne sert qu'a
%  cela. Le tableau 11 romanche avait montre deux petites langues
%  faisant le meme partage sur la caisse d'epargne et les pompiers, et
%  concluait que la cause n'etait dans aucune des deux. Ici c'est le
%  meme raisonnement sur un mois.
%
%  MAIS LA DIFFERENCE COMPTE AUTANT QUE L'ACCORD. Le romanche n'a
%  refait QUE DEUX mois : les dix autres sont restes latins.
%  L'estonien a nomme LES DOUZE, et ses douze se partagent en deux
%  moities egales, six travaux et six saints. Le romanche montre la
%  ferme et Rome ; l'estonien montre la ferme et l'eglise, et Rome
%  n'y est pas. La ou une langue a garde un calendrier entier, ce
%  n'est pas le travail seul qui l'a fait, c'est le travail et la
%  fete.
%
%  LA COLONNE ECRIT LES MOIS LATINS. La planche est un manuel
%  scolaire, l'ecole estonienne date en « märts, aprill, mai », et
%  c'est la forme que le lecteur cherchera. Les noms paysans sont
%  releves ici, ou ils s'expliquent, et non dans le corps, ou ils
%  feraient un contresens de registre.
%
%  ECART DE COMPOSITION DU BLOC t03-c2-01-1 : le fac-simile ido y
%  ecrit « junio (julio od agosto) » avec « julio » en romain, seul
%  des trois. gravuri/korekti.json remet le gras ; le fichier source
%  ne bouge pas.
%
%  UNE INVERSION, LA PREMIERE DE LA COLONNE, AU BLOC t03-c1-04-1.
%  L'ido range hirondelles (4), fleche (7), clocher (5), eglise (6) ;
%  l'estonien enchaine ses genitifs a l'envers du francais — « küla
%  kiriku kellatorni tornikiiver », le village de l'eglise du clocher
%  la fleche — et la premiere redaction donnait donc 4, 6, 5, 7.
%  renvoji.py l'a signalee. La phrase a ete coupee en deux : la
%  fleche d'abord avec son clocher, puis une seconde proposition pour
%  l'eglise que le clocher couronne. La chaine de genitifs disparait
%  et l'ordre de l'ido revient sans qu'on force la langue.
% ===================================================================

%%K t03-apar-1 apar
\VUsaut{3.0mm}
\VUtitre{15.0pt}{60}{TABEL Nr. 3}
\VUsaut{1.6mm}
\VUtitre{11.4pt}{0}{Lapsepõlv ja noorus.}
\VUsaut{3.4mm}

%%K t03-apar-2 apar
(Tabelid 3 ja 4 on kokku seatud nii, et nad esitavad neljas \VUgras{perekonnastseenis}
peamised mõisted \VUgras{ilma}, \VUgras{vanuse}, \VUgras{perekonna},
\VUgras{riietuse} ja \VUgras{seisundi} kohta.)

\VUsaut{4.0mm}
%%K t03-tit-1 sub
\VUcentre{12.0pt}{0}{\textit{Esimene stseen.}}
\VUsaut{3.0mm}

%%K t03-tit-2 sub
\VUpk{10.2pt}{40}{Ristimistalitus külas.}
\VUsaut{2.4mm}

%%K t03-tit-3 sub
\VUpk{10.2pt}{40}{Kevad.}
\VUsaut{3.0mm}

%%K t03-c1-01-1 p
1. --- Me oleme \VUgras{kevadel}, kuudel \VUgras{märts}, \VUgras{aprill} või
\VUgras{mai}, nädalapäevadel \VUgras{esmaspäev}, \VUgras{teisipäev} või
\VUgras{kolmapäev}. \VUgras{Kell on kümme} hommikul.

%%K t03-c1-02-1 p
2. --- \VUgras{Ilm} on ilus, \VUgras{õhk} puhas. Päike paistab. Mõned
\VUgras{pilvekesed} \textsuperscript{(2)} tõusevad sinises \VUgras{taevas}
\textsuperscript{(1)}.

%%K t03-c1-03-1 p
3. --- \VUgras{Puudel} on vaevalt \VUgras{lehti}. Saledatel \VUgras{paplitel}
\textsuperscript{(3)} ja kõrgel \VUgras{plataanil} \textsuperscript{(17)} on seni ainult
pungad.

%%K t03-c1-04-1 p
4. --- \VUgras{Pääsukesed} \textsuperscript{(4)} tulevad tagasi ja lendlevad rõõmsa
siutsumisega ümber \VUgras{tornikiivri} \textsuperscript{(7)}, mis on
\VUgras{kellatorni} \textsuperscript{(5)} otsas; see torn kroonib küla
\VUgras{kirikut} \textsuperscript{(6)}.

%%K t03-tit-4 sub
\VUsaut{3.4mm}
\VUpk{10.2pt}{40}{Kirik.}
\VUsaut{3.0mm}

%%K t03-c1-05-1 p
5. --- Väga lihtne on see kirik oma suure \VUgras{peauksega} \textsuperscript{(13)} ja
suure \VUgras{kellaga} \textsuperscript{(8)}. Väike \VUgras{osuti} \textsuperscript{(9)}
on numbril X, see näitab \VUgras{tunde}. \VUgras{Suur} \textsuperscript{(10)} on numbril
XII (keskpäev); see näitab \VUgras{minuteid}.

%%K t03-c1-06-1 p
6. --- Kellatornis helisevad \VUgras{kellad} \textsuperscript{(11)} rõõmsalt. Katuse
harjal seisab väike kivist \VUgras{rist} \textsuperscript{(12)}.

%%K t03-c1-07-1 p
7. --- Kirikul ei ole ainult suur \VUgras{peauks} \textsuperscript{(13)}, tal on ka
väike külguks. Sellest uksest väljub härra \VUgras{kirikuõpetaja}
\textsuperscript{(15)}, riietatud oma \VUgras{sutaani}, \VUgras{käärkambrist}
\textsuperscript{(14)} ja läheb \VUgras{pastoraati} \textsuperscript{(19)}.

%%K t03-c1-08-1 p
8. --- Tema lähedal on \VUgras{piimanaine} \textsuperscript{(24)} oma
\VUgras{eesliga} \textsuperscript{(23)}. Ta tuleb tagasi linnast, kuhu ta viis head
piima laste jaoks.

%%K t03-tit-5 sub
\VUsaut{4.0mm}
\VUpk{10.2pt}{40}{Viljapuuaed.}
\VUsaut{3.4mm}

%%K t03-c1-09-1 p
9. --- Härra Aliboron (eesel) on selle hommikuse jooksuga täiesti rahul. Ta hingab
mõnuga sisse õhku, mida on lõhnastanud \VUgras{õied} \textsuperscript{(20)}, mis katavad
naabruses oleva \VUgras{viljapuuaia} \VUgras{viljapuid} \textsuperscript{(18)}. Täiesti
roosad ja valged on need \VUgras{virsikupuud} \textsuperscript{(18)} ja
\VUgras{aprikoosipuud} \textsuperscript{(19)}, ja nad lubavad head saaki.

%%K t03-c1-10-1 p
10. --- Vahepeal lähevad väikesed \VUgras{mesilased} \textsuperscript{(22)}
\VUgras{mesipuust} \textsuperscript{(21)} sinna suminal oma magusat \VUgras{mett}
korjama.

%%K t03-c1-11-1 p
11. --- Kuid mis juhtub? Vaat et küla poisid jooksevad kokku ja ajavad hirmunud
\VUgras{haned} \textsuperscript{(25)} laiali. See on rongkäigu väljumine kirikust pärast
notari poja \VUgras{ristimist}.

%%K t03-c1-12-1 p
12. --- Väike Jaakob ärkab oma \VUgras{hällis} \textsuperscript{(26)} kellade rõõmsast
helinast, ja tema õde Magdalena, kes tahab samuti ristimisrongkäiku näha, palub oma
ema tulla tema juukseid kammima ja teda riietama maja lävel.

%%K t03-c1-13-1 p
13. --- Ema nõustub. Ta paneb pähe oma \VUgras{rätiku} \textsuperscript{(28)}, selga oma
\VUgras{pihiku} \textsuperscript{(29)} ja ette oma \VUgras{põlle} \textsuperscript{(30)},
ja tuleb istuma väiksele toolile ukse ette. Magdalenal on ainult tema
\VUgras{alusseelik} \textsuperscript{(31)} ja tema \VUgras{korsett}
\textsuperscript{(32)}.

%%K t03-c1-14-1 p
14. --- Kui õnnelik ta on! Ta unustab oma lemmiklilled, oma \VUgras{nelgid}
\textsuperscript{(34)}, millele laskuvad \VUgras{liblikad} \textsuperscript{(35)}, oma
\VUgras{tulbid} \textsuperscript{(36)}, oma \VUgras{maikellukesed} \textsuperscript{(37)}
\VUgras{pottides} \textsuperscript{(38)} \VUgras{müüril} \textsuperscript{(33)}, ja oma
\VUgras{roosid} \textsuperscript{(39)}, mis moodustavad nii ilusa heki. Ta vaatab
saatjaskonna daamide rikkalikke riideid.

%%K t03-c1-15-1 p
15. --- Küla poisid ei pane riietust eriti tähele. Nad tormavad ja tõuklevad, et
saada \VUgras{kompvekke} \textsuperscript{(55)} või \VUgras{münte} \textsuperscript{(52)}.
Üks neist nutab \textsuperscript{(40)}.

%%K t03-c1-16-1 p
16. --- Teine astus \VUgras{koera} \textsuperscript{(42)} käpa peale, kes kiunudes
põgeneb ja ajab hirmu \VUgras{kanale} \textsuperscript{(43)} ja tema \VUgras{tibudele}
\textsuperscript{(44)}.

%%K t03-tit-6 sub
\VUsaut{5.0mm}
\VUpk{10.2pt}{40}{Ristimisrongkäik.}
\VUsaut{4.0mm}

%%K t03-c1-17-1 p
17. --- Rongkäigu ees kõnnib \VUgras{amm} \textsuperscript{(45)}. Tal on ilus
\VUgras{tanu} \textsuperscript{(46)} pikkade paeltega. Ta kannab \VUgras{vastsündinut}
\textsuperscript{(47)}, kes on üleni \VUgras{pitsidesse} \textsuperscript{(48)}
riietatud.

%%K t03-c1-18-1 p
18. --- Amme järel tulevad \VUgras{ristiisa} \textsuperscript{(49)} ja
\VUgras{ristiema} \textsuperscript{(53)}. Nad jagavad kompvekke ja münte. Ristiisal on
\VUgras{kindad} \textsuperscript{(51)} ja \VUgras{silinderkübar} \textsuperscript{(50)}.
Ristiema hoiab käes ilusat \VUgras{lillekimpu} \textsuperscript{(54)}.

%%K t03-c1-19-1 p
19. --- Nende järel näeme me lapse \VUgras{onu} \textsuperscript{(56)} ja
\VUgras{tädi} \textsuperscript{(58)}. Tädil on elegantne \VUgras{kübar}
\textsuperscript{(59)}, onul must \VUgras{redingot} \textsuperscript{(57)} ja valge vest.
Siin on siis \VUgras{onupoeg} \textsuperscript{(60)} ja \VUgras{onutütar}
\textsuperscript{(61)}. Viimane hoiab käest vastsündinu \VUgras{õde}
\textsuperscript{(62)}, väikest Bertat.

%%K t03-c1-20-1 p
20. --- Berta teine \VUgras{vend} \textsuperscript{(63)} kõnnib taga. Ta annab käe ühele
\VUgras{sugulasele} \textsuperscript{(64)}. Perekonna \VUgras{sõbrad}
\textsuperscript{(65)} tulevad siis.

%%K t03-tit-7 sub
\VUsaut{4.6mm}
\VUpk{10.2pt}{40}{Pealtvaatajad.}
\VUsaut{3.4mm}

%%K t03-c1-21-1 p
21. --- Rongkäigu lähedal on \VUgras{kerjus} \textsuperscript{(66)}, kes toetub oma
\VUgras{kargule} \textsuperscript{(67)}. Ta palub almust.

%%K t03-c1-22-1 p
22. --- Teisel pool on väike väga \VUgras{viisakas} \textsuperscript{(68)} poiss. Ta
tervitab ja soovib «\VUgras{tere päevast}» inimestele, keda ta tunneb.

%%K t03-c1-23-1 p
23. --- Külaelanikud on oma majadest välja tulnud ristimisrongkäiku vaatama. Me näeme
\VUgras{talumeest} \textsuperscript{(69)} tema \VUgras{puuvillase mütsiga} ja tema kõrval
\VUgras{talunaist} \textsuperscript{(70)}. Siis \VUgras{vanaeite} \textsuperscript{(71)},
kes toetub oma \VUgras{kepile} \textsuperscript{(72)}. Lõpuks \VUgras{möldrit}
\textsuperscript{(73)} tema suure \VUgras{viltkübaraga} \textsuperscript{(74)} ja noort
talunaist \VUgras{puukingades} \textsuperscript{(75)}, kes õpetab oma väikelast käima.

%%K t03-c1-24-1 p
24. --- See stseen on väga lõbus.

\VUsaut{4.0mm}
%%K t03-tit-8 sub
\VUcentre{12.0pt}{0}{\textit{Teine stseen.}}
\VUsaut{3.0mm}

%%K t03-tit-9 sub
\VUpk{10.2pt}{40}{Avalik pidu.}
\VUsaut{2.4mm}

%%K t03-tit-10 sub
\VUpk{10.2pt}{40}{Veemängud ja regatid.}
\VUsaut{3.0mm}

%%K t03-c2-01-1 p
1. --- \VUgras{Suvi} on kätte jõudnud. \VUgras{Juuni} (juuli või \VUgras{august}) kuu
palav päike ajab noored suplema ja paadiga sõitma.

%%K t03-c2-02-1 p
2. --- Jõe paremal kaldal võtavad paadimehed riided seljast, et võtta osa veemängudest
ja regatidest.

%%K t03-c2-03-1 p
3. --- Üks neist võtab jalast oma \VUgras{kingad} \textsuperscript{(1)}; ta nööpib need
lahti. Tema kaks kaaslast on juba valmis. Nüüd on neil ainult oma \VUgras{trikoosärk}
\textsuperscript{(2)} ja \VUgras{ujumispüksid} \textsuperscript{(3)}. Nad vestlevad, oma
sõpru oodates. Nad räägivad laadast ja peost, mis toimub teisel kaldal.

%%K t03-c2-04-1 p
4. --- Maas, nende kõrval, lebavad nende kaaslaste \VUgras{rõivad}: \VUgras{paar}
\VUgras{poolsaapaid} \textsuperscript{(4)}, \VUgras{kingi} \textsuperscript{(5)},
\VUgras{sokke} \textsuperscript{(6)}, \VUgras{vildist} \textsuperscript{(7)} ja
\VUgras{õlgedest} \textsuperscript{(8)} kübaraid ning \VUgras{lips}
\textsuperscript{(9)}.

%%K t03-c2-05-1 p
5. --- Üks noormeestest on hoolikalt kokku voltinud oma \VUgras{ülikonna}
\textsuperscript{(10)}. Ta paneb selle käterätikule, millesse tema naaber peagi mähib
mõlemad \VUgras{rõivad}.

%%K t03-c2-06-1 p
6. --- Kuues poiss jääb hiljaks. Tal on veel \VUgras{nokkmüts} \textsuperscript{(11)}
peas. Ta ei ole seljast võtnud ei oma \VUgras{särki} \textsuperscript{(12)}, ei oma
\VUgras{kraed} \textsuperscript{(13)}, ei oma \VUgras{kätiseid} \textsuperscript{(14)}. Ta
hoiab käes oma \VUgras{vesti} \textsuperscript{(17)} ja tal on veel jalas oma
\VUgras{püksid} \textsuperscript{(16)} ja \VUgras{traksid} \textsuperscript{(15)}. Kindlasti
ei ole ta järgmiseks võiduajamiseks valmis.

%%K t03-c2-07-1 p
7. --- Noormeeste kõrval viib rada, mille servadel on palju \VUgras{ohakaid}
\textsuperscript{(18)}, jõe kaldale, kus isa Jaakob valmistab \VUgras{pilliroo}
\textsuperscript{(19)} vahel \VUgras{ilutulestikku} \textsuperscript{(20)}, mida õhtul
lastakse.

%%K t03-tit-11 sub
\VUsaut{4.4mm}
\VUpk{10.2pt}{40}{Võiduajamine.}
\VUsaut{3.4mm}

%%K t03-c2-08-1 p
8. --- Jõel jalutavad mõned daamid \VUgras{paadis} \textsuperscript{(21)}, varjates end
oma päikesevarjudega. Nad jälgivad võiduajamist, mis on väga huvitav. Igas
\VUgras{joolis} \textsuperscript{(22)} sõuavad neli \VUgras{sõudjat}
\textsuperscript{(24)} kõigest jõust, sellal kui \VUgras{tüürimees}
\textsuperscript{(26)} hoiab \VUgras{tüüri} \textsuperscript{(25)} ja juhib haprat
sõudepaati. \VUgras{Aerud} \textsuperscript{(23)} löövad taktis vett ja kerged paadid
liuglevad peadpööritava kiirusega.

%%K t03-c2-09-1 p
9. --- Mõned noormehed võistlevad silla samba juures \VUgras{libeda masti}
\textsuperscript{(28)} auhinna pärast. Üks neist astub ettevaatlikult mööda painduvat
ja libisevat masti, et ulatuda väikese \VUgras{lipukeseni} \textsuperscript{(29)}, mis on
otsa kinnitatud. Tema õnnetu \VUgras{võistleja} \textsuperscript{(30)} on vette
kukkunud. Ta ujub, et kaldale saada.

%%K t03-c2-10-1 p
10. --- Vanemad noormehed peavad \VUgras{paaditurniiri} \textsuperscript{(32)}
\VUgras{kai} \textsuperscript{(33)} lähedal. Kõiki neid mänge vaatab arvukas
\VUgras{publik} \textsuperscript{(31)}, kes toetub \VUgras{silla} \textsuperscript{(27)}
\VUgras{käsipuule} ja on kogunenud \VUgras{kaldale}. Mõned pealtvaatajad pöörduvad
siiski ümber, et jälgida \VUgras{auhinnamasti} \textsuperscript{(45)} mängu või imetleda
\VUgras{linnavalitsuse rongkäiku}, mis tuleb auhindade \VUgras{jagamisele}.

%%K t03-c2-11-1 p
11. --- Kõigepealt on siin mõned \VUgras{poisid} \textsuperscript{(35)}, kes käivad
\VUgras{muusikute} \textsuperscript{(36)} ees; siis tulevad \VUgras{linnapea}
\textsuperscript{(37)}, abilised ja linnakese \VUgras{linnavolikogu}. Lõpuks marsivad
\VUgras{tuletõrjujad} \textsuperscript{(38)}.

%%K t03-tit-12 sub
\VUsaut{5.0mm}
\VUpk{10.2pt}{40}{Laat.}
\VUsaut{3.6mm}

%%K t03-c2-12-1 p
12. --- \VUgras{Laadaplatsil} \textsuperscript{(39)} on arvukalt rahvast. Tullakse
vaatama mitmesuguseid \VUgras{putkasid}: \VUgras{puuhobuseid} \textsuperscript{(40)},
\VUgras{tsirkust} \textsuperscript{(41)}, kus etendus on juba alanud; \VUgras{avalikku
balli} \textsuperscript{(42)}, kus linna noormehed ja neiud naudivad
\VUgras{tantsimise} rõõmu.

%%K t03-c2-13-1 p
13. --- Tagaplaanil näeme me \VUgras{areeni} \textsuperscript{(44)}, kus vapper hispaania
\VUgras{matadoor} võitleb raevunud \VUgras{härjaga} \textsuperscript{(45)}, siis
\VUgras{ameerika mägesid} \textsuperscript{(49)}, \VUgras{lasketiiru}
\textsuperscript{(46)}, kus harjutatakse \VUgras{tulirelvade} kasutamist ja
\VUgras{märklauda} laskmist.

%%K t03-c2-14-1 p
14. --- Lähedal on \VUgras{laadateater} \textsuperscript{(47)}, kus mängitakse
\VUgras{komöödiat} ja \VUgras{draamat}. Väga lähedal on \VUgras{hiiglase}
\textsuperscript{(48)} ja \VUgras{kääbuse} putka. Esimese \VUgras{kasv} on kaks meetrit
ja viisteist sentimeetrit; teine on vaevalt nii suur kui laps.

%%K t03-c2-15-1 p
15. --- Lõpuks on siin suur \VUgras{loomaaed} \textsuperscript{(50)} oma
\VUgras{metsloomadega}, oma \VUgras{tiigriga} \textsuperscript{(51)}, kellel on võimsad
\VUgras{küüned}, oma \VUgras{lõviga} \textsuperscript{(52)}, kellel on tihe \VUgras{lakk},
oma kahe \VUgras{kaelkirjakuga} \textsuperscript{(53)} ja oma \VUgras{elevandiga}
\textsuperscript{(54)}, kes nii osavalt oma \VUgras{lonti} kasutab.

%%K t03-c2-16-1 p
16. --- \VUgras{Taltsutaja} \textsuperscript{(55)}, kes neid loomi treenib, on putka ukse
juures.

\VUsaut{6.0mm}
\VUfilet{22mm}
